\scanpage{20}
\mspage{18}
% archive sheet number: 19
% hand: Grothendieck
% sheet_kind: text
% element: text | top, line 1 | En d'autres termes
En d'autres termes:

% element: text | upper, lines 2–5 | Corollaire 4.5
\uline{Corollaire} 4.5. Avec les notations de \unsure{4.4}, quitte à \deleted{\unsure{par}} la restreindre à un sous-groupe ouvert \inserted{convenable} $U$ de \unsure{$\pi$|$I$}, la représentation donnée $u:I\longrightarrow G$ se factorise en une représentation
% element: display | upper, line 6 | U to Im(U to Z_l(1)) to G
\[
U\longrightarrow\operatorname{Im}(U\longrightarrow\mathbb Z_\ell(1))\unsure{-|\longrightarrow}G
\]
% element: text | upper-middle, lines 7–8 | d'un sous-groupe ouvert
d'un sous-groupe ouvert \inserted{$U'$\deleted{\unsure{$U$}}} du facteur \deleted{\unsure{$\mathbb Z_\ell(1)$}} $\ell$-primaire $\mathbb Z_\ell(1)$ de $I_t$.

% element: text | middle, lines 9–10 | Si l^nu est l'indice
Si $\ell^\nu$ est l'indice de ce dernier, alors la multiplication par $\ell^\nu$ induit un isomorphisme
% element: display | middle, line 11 | Z_l(1) isom U'
\[
\mathbb Z_\ell(1)\xrightarrow{\sim}U',
\tag*{(4.7)}
\]
% element: text | middle, lines 12–14 | et on est donc ramené
et on est donc ramené, pour étudier la restriction de $u$ à $U$, à étudier une représentation $\ell$-adique du groupe $\mathbb Z_\ell(1)$.

% element: text | lower-middle, lines 15–16 | En fait, souvent
En fait, souvent on se borne à prendre $U$ tel que
% element: display | lower-middle, lines 16–17 | P cap U subset P cap N
\[
\begin{aligned}
&\deleted{\unsure{K\subset N}\quad K\cap U\subset\unsure{K}}\\
&P\cap U\subset P\cap N,
\end{aligned}
\]
% element: text | lower, lines 18–19 | de sorte que la restriction
de sorte que la restriction de $u$ à $U$ se factorise en un hom
% element: display | lower, line 20 | U' = Im(U to I_t) to G
\[
U'=\operatorname{Im}(U\longrightarrow I_t)\longrightarrow G.
\]
% element: text | lower, line 21 | Ici encore
Ici encore, on trouve un isom canonique
% element: display | lower, line 22 | U' isom I_t isom product
\[
U'\simeq I_t\simeq\prod_{\ell'\ne p}\mathbb Z_{\ell'}(1)
\tag*{(4.8)}
\]
% element: text | bottom, line 23 | On retiendra
On retiendra en particulier que $U'$ est \uline{commutatif}.

% element: text | bottom edge, lines 24–27 | 4.9 On peut aussi exprimer
\inserted{4.9.} On peut aussi exprimer ceci en disant que, quitte à faire une extension séparable finie $K_1$ de $K$, la repr. $\ell$-adique donnée de $I$ définit une repr. $\ell$-adique \unsure{de $I_1$} qui se factorise par sa ``partie $\ell$-primaire'' $\unsure{I(\ell)|I_1(\ell)}=\mathbb Z_\ell(1)$.
